\input{\string~/Documents/School/"UC Irvine"/ta_template2.0.tex}
\rh{2B}{11.1}
\chead{\today}
\graphicspath{ {\string~/Desktop} }
\usetikzlibrary{arrows}
%============ANSWER KEY TOGGLE===========
\newcommand{\ans}[1]{ \noindent {\color{red} \textbf{Answer:} #1}}
\renewcommand{\ans}[1]{\vspace{3in}}
%==========================================
\begin{document}
\begin{center}
\textbf{Sequences}
\end{center}

\begin{defi}[Limit of a Sequence]
A sequence (think list of numbers $a_1, a_2, \cdots$) has a limit $L$ if for every $\ee > 0$, there exists an integer $N$ such that for all $n > N$, $\lrabs{a_n - L} < \ee$.
\end{defi}
\begin{thm}
If $\lim_{x \to \infty} f(x) = L$ and $f(n) = a_n$ for some sequence $a_n$, then $\lim_{n \to \infty} a_n = L$. 
\label{limitfn}
\end{thm}
\begin{aun}
Remember that limits preserve sums ($\lim (a+b) = \lim a + \lim b$), differences, products, quotients, scalar multiples, and powers (in a sense they just work really nicely for us). 
\end{aun}
\begin{thm}[Squeeze Theorem] If $a_n \leq b_n \leq c_n$ and we have
\[
L = \lim_{n \to \infty} a_n = \lim_{n \to \infty} c_n \mte{ then } \lim_{n \to \infty} b_n  = L
\]
\end{thm}
\begin{thm}
If $\lim \lrabs{a_n} = 0$ then $\lim a_n  = 0$.
\label{conv}
\end{thm}
\begin{thm}
If $\lim_{n \to \infty} a_n = L$ and the function $f$ is continuous at $L$, we then:
\[
\lim_{n \to \infty} f(a_n) = f(L)
\]
\end{thm}
\begin{thm}[Geometric Sequences]
The sequence $(r^n)_{n \in \N}$ is convergent if $-1 < r \leq 1$ and diverges for all other values:
\[
\lim _{n \to \infty} r^n = \begin{cases}
0 & \lrabs{r} < 1\\
1 & r = 1
\end{cases}
\]
\end{thm}
\begin{aun}
I won't define bounded or monotonic really precisely here, but think ``bounded below'' as just a number that's lower than every point of the sequence, ``bounded above'' is a number above every point of the sequence, and monotonic just means the sequence is always increasing or always decreasing.
\end{aun}
\begin{aun}[Monotone Sequence Theorem] If we can show that a sequence is bounded and monotonic, then it converges.
\end{aun}
\begin{act}
In groups of up to 3, draw pictures to help your intuition on \textbf{Theorems 5, 6, 8} and \textbf{Author's Note 4}. I will come around and discuss these with you!\\
Once your group is finished, draw a picture of the sequence:
\[
a_1 = 2, \mte{ and } a_{n+1} = \frac{a_n}{1+ a_n}
\]
\end{act}

\ans{\href{https://www.desmos.com/calculator/5rtmqiqpv2}{Use this link}}

\problembox{Find a formula for the following sequences, assuming the pattern continues:
\begin{enumerate}[label=\alph*.]
\item $\lrsq{\frac12, \frac14, \frac{1}{6}, \frac{1}{8}, \frac{1}{10}, \cdots}$
\item $\lrsq{1, 1, 2, 3, 5, 8, 13, 21, \cdots}$
\item $\lrsq{\frac12, \frac{-4}{3}, \frac94, \frac{-16}{5}, \frac{25}{6}, \cdots}$
\end{enumerate}
}

\ans{
\begin{enumerate}[label=\alph*.]
\item $a_n = \frac{1}{2n}$
\item $a_{n} = a_{n - 1} + a_{n-2}$
\item $\dst \frac{(-1)^{n+1} n^2}{n+1}$
\end{enumerate}
}

\problembox{Determine whether the following sequences converge:
\[
\lrsq{\sin(n)} \mte{ and } a_n = \frac{(-1)^n}{2 \sqrt{n}}
\]}

\ans{
First doesn't, it doesn't get ``close'' to anything.\\
Second will converge. Using Theorem~\ref{conv}, we see that we can evaluate the limit of the absolute value first, which is:
\[
\lrabs{a_n} = \frac{1}{2 \sqrt{n}}
\]
Using Theorem~\ref{limitfn}, it suffices to check the end behavior of the function
\[
\lim_{x \to \infty} \frac{1}{2\sqrt{x}} = 0
\]
meaning that the sequence $(a_n)$ converges as well. 
}

\problembox{Consider the sequence:
\[
a_n = \arctan \lrp{ \ln (n)}
\]
Does it converge? If so, what to?
}

\ans{Yes, to $\frac{\pi}{2}$; use Theorem~\ref{limitfn}}

\end{document} 